\homework{8}{Sun 23 Oct 2022 14:00}{}

Herstein, section $2.10$, page 101 , Questions 1,3 .
\begin{itemize}
	\item 1. Let $A$ be a normal subgroup of a group $G$, and suppose that $b \in G$ is an element of prime order $p$, and that $b \notin A$. Show that $A \cap(b)=(e)$.
\begin{proof}
	Suppose $A \cap \left\langle b \right\rangle $ is nontrivial, so we can pick $b^i \in A \cap \left\langle b \right\rangle $ where $b^i \neq  e$ and $b^i \in A$. Now we want to prove $b \in A$.

	Since $A \triangleleft G$ and $b^i \in A$, for all $g \in G$ we have $g^{-1} b^i g \in A$. Since $o(b) = p$ and $p$ prime, we have $(i, p) = 1$. By Euclidean Algorithm, there exists integers $k, m$ such that $ki + mp = 1$, so $ki \equiv 1 \pmod p$. Therefore
	\[
		\left( g^{-1} b^i g \right)^k =  g^{-1} b^{ik} g = g^{-1} bg \in A
	.\] 
	Hence $bg \in  gA$. But $gA = Ag$, so $bg \in Ag$. This implies $b \in A$.
\end{proof}
\item
3. Let $G$ be a finite group, with $N \triangleleft G$ and $a \in G$. Prove that:\\
(a) The order of $a N$ in $G / N$ divides the order of $a$ in $G$, that is, $o(a N) \mid o(a)$.
\begin{proof}
	We have
	\[
		(aN)^{o(a)} = a^{o(a)}N = eN = N
	.\] 
	This $o(aN) | o(a)$. This follows from a past exercise saying that if $a^s = e$ then $o(a) | s$.
\end{proof}

(b) If $(a) \cap N=(e)$, then $o(a N)=o(a)$.
\begin{proof}
	Since $(aN)^{o(aN)} = a^{o(aN)} N= N$, we have $a^{o(aN)} \in \left\langle a \right\rangle \cap N = (e)$. Thus $o(a) | o(aN)$. Thus $o(a) = o(aN)$.
	
\end{proof}
\item
Additional question (2.10.A): Give a representative for every isomorphism class of abelian groups of order 240.
\begin{proof}
	Factorize: $240 = 2^4 3^1 5^1$.
	The representatives are
	\begin{itemize}
		\item $\mathbb{Z}_{30}\times \mathbb{Z}_{2}\times \mathbb{Z}_2 \times \mathbb{Z}_2$
		\item $\mathbb{Z}_{60}\times \mathbb{Z}_2 \times \mathbb{Z}_2$
		\item $\mathbb{Z}_{60}\times \mathbb{Z}_4 $
		\item $\mathbb{Z}_{120}\times \mathbb{Z}_2 $
		\item $\mathbb{Z}_{240}$
	\end{itemize}
\end{proof}
\item
Additional question (2.10.B): Give a representative for every isomorphism class of abelian groups of order 540 .
In each case, make sure that your answer is a direct product involving no more than four factors. 
\begin{proof}
	Factorize: $540 = 2^2 3^3 5^1$
	The representatives are
	\begin{itemize}
		\item $\mathbb{Z}_{30} \times \mathbb{Z}_3 \times \mathbb{Z}_3 \times \mathbb{Z}_2 \cong \mathbb{Z}_{30} \times \mathbb{Z}_6 \times \mathbb{Z}_3$
		\item $\mathbb{Z}_{90} \times \mathbb{Z}_3 \times \mathbb{Z}_2 \cong \mathbb{Z}_{90} \times \mathbb{Z}_6$
		\item $\mathbb{Z}_{270} \times \mathbb{Z}_2$
		\item $\mathbb{Z}_{60} \times \mathbb{Z}_3 \times \mathbb{Z}_3 $
		\item $\mathbb{Z}_{180} \times \mathbb{Z}_3 $
		\item $\mathbb{Z}_{540}$
	\end{itemize}
\end{proof}


\end{itemize}
\newpage
Herstein, section 2.11, page 106, Questions 1, 3, 4, 5, 9, 13.
\begin{itemize}[resume]
	\item 1. In $S_3$, the symmetric group of degree 3 , find all the conjugacy classes, and check the validity of the class equation by determining the orders of the centralizers of the elements of $S_3$.
\begin{proof}
	$S_3 = \{e, (1,2),(1,3), (2,3), (1,2,3), (3,2,1)\} $
	 \begin{itemize}
		 \item $\operatorname{cl}(e) = \{e\} $
		 \item $\operatorname{cl}((1,2)) = \{(1,2), (2,3), (1,3)\} $
		 \item $\operatorname{cl}((1,2,3)) = \{((1,2,3), (3,2,1)\} $
	 \end{itemize}
	 Also,
	 \begin{itemize}
	 	\item $C(e) = S_3$, so $i_{S_3}(C(e)) = 1$
	 	\item $C((1,2)) = \{e, (1,2)\} $, so $i_{S_3}(C(e)) = \frac{6}{2} = 3$
	 	\item $C((1,2,3)) = \{e, (1,2,3),(3,2,1)\} $, so $i_{S_3}(C(e)) = \frac{6}{3} = 2$
	 \end{itemize}
	 We verify that
	 \begin{itemize}
	 	\item $\left| \operatorname{cl}(e) \right|  =  1 = i_{S_3}(C(e))$
	 	\item $\left| \operatorname{cl}((1,2)) \right|  =  3 = i_{S_3}(C((1,2)))$
	 	\item $\left| \operatorname{cl}((1,2,3)) \right|  = 2  = i_{S_3}(C((1,2,3)))$
	 \end{itemize}
	 Since $1 + 3 + 2 = 6 = 3! = |S_3|$, we have verified the class equation.
\end{proof}

	\item 3. If $a \in G$, show that $C\left(x^{-1} a x\right)=x^{-1} C(a) x$.
\begin{proof}
	Take any $g \in G$.
	\begin{align*}
		g \in C(x^{-1} ax) \\
		&\iff g x^{-1} a x = x^{-1} a x g \\
		&\iff g^{-1} x^{-1} a x g = x^{-1}  a x \\
		&\iff x^{-1} a x g x^{-1}  a^{-1} x = g \\
		&\iff a  x g x^{-1} a^{-1} = x g x^{-1} \\
		&\iff a x g x^{-1}  = x g x^{-1} a \\
		&\iff  x g x^{-1}  \in C(a) \\
		&\iff g \in x^{-1}  C(a) x 
	.\end{align*}
\end{proof}

	\item 4. If $\varphi$ is an automorphism of $G$, show that $C(\varphi(a))=\varphi(C(a))$ for $a \in G$.
\begin{proof}
	Take any $g \in G$.  Now
	\begin{align*}
		g \in C(\varphi(a))\\
		&\iff g \varphi(a) = \varphi(a) g \\
		&\iff \varphi(\varphi^{-1} (g))\varphi(a) = \varphi(a) \varphi(\varphi^{-1} (g)) & \text{since $\varphi$ is bijective} \\
		&\iff \varphi(\varphi^{-1} (g) a) = \varphi(a \varphi^{-1} (g)) &\text{since $\varphi$ is homomorphism} \\
		&\iff \varphi^{-1} (g) a = a \varphi^{-1} (g) &\text{since $\varphi$ is bijective} \\
		&\iff \varphi^{-1} (g) \in C(a) \\
		&\iff g \in \varphi(C(a)) 
	.\end{align*}
\end{proof}

	\item 5. If $|G|=p^3$ and $|Z(G)| \geq p^2$, prove that $G$ is abelian.
\begin{proof}
	We assume $p$ is prime. Now $|G / Z(G)| \le p$, so $|G / Z(G)| = 1 $ or $p$. If  $|G / Z(G)| = 1$ then $|Z(G)| = |G|$, so $Z(G) = G$, hence $G $ abelian.
	If $|G / Z(G)| = p$, then  $Z(G) = \left( e \right) $. However this cannot happen since by theorem 2.11.4, the center of $G$ is nontrivial.

If $p$ is not prime, this statement is not always true. We show this by giving a counterexample. Let $A$ be an abelian group of order $p^2$ and $B$ be a non-abelian group of order $p$ (such group can sometimes exist since $p$ is not prime). Then the group $A \times B$ is clearly nonabelian; but it has order $p^3$ and $|Z(A \times B)| \ge  |A| \times |Z(B)| \ge  |A| = p^2$. 
\end{proof}

	\item 9. Prove that $N(H)$ is a subgroup of $G, H \subset N(H)$ and in fact $H \triangleleft N(H)$.
\begin{proof}
	Let $a, b \in N(H)$. We want to show $ab^{-1}  \in N(H)$.
	\begin{align*}
		a,b \in N(H)\\
		&\iff aH = Ha, bH = Hb \\
		&\implies a H = a b^{-1} b H = ab^{-1} H b \\
		&\iff ab^{-1} H = aH b^{-1}  = H ab^{-1}  \\
		&= ab^{-1}  \in N(H)
	.\end{align*}
	Hence by subgroup criterion, $N(H) \le  G$.

	To see $H \subset N(H)$, note that for any  $h \in H$, $hH = H = Hh$.

	To see that $H \triangleleft  N(H)$, take any element $n \in N(H)$, and we want to show $n H = H n$. But this directly follows from definition of $N(H)$.
\end{proof}

	\item 13. Prove that if $G$ is a finite group and $H$ is a subgroup of $G$, then the number of distinct subgroups $x^{-1} H x$ of $G$ equals $i_G(N(H))$.
\begin{proof}
	Define $\varphi$ to be the mapping $x^{-1} H x \mapsto N(H) x$.
	Observe that for any $x, y \in G$,
	\begin{align*}
		x^{-1} H x  = y^{-1} H y \\
		&\iff (xy^{-1} )^{-1} H (xy^{-1} ) = H \\
		&\iff xy^{-1} \in N(H) \\
		&\iff N(H) x = N(H)y
	.\end{align*}
	So $\varphi$ is a bijection from all subgroups $x^{-1} H x$ to $G / N(H)$. So the number of distinct $x^{-1} H x$ is exactly equal to $|G / N(H)| = i_G(N(H))$.
\end{proof}

\end{itemize}
